\section{供用电、水、气热力合同、赠与合同}
\begin{enumerate}
  \item 不是考试重点。
  \item 供用电、水、气热力合同。\al{经催告吴某在合理期限内未交费，可以解除供热合同 ×}。
        \begin{enumerate}
          \item 强制缔约 \yjt 不可以解除。\al{经催告吴某在合理期限内未交费，\xhx{可以中止}供热}。
          \item 免责事由仅限\zd{不可抗力 + 及时采取抢修措施}。
        \end{enumerate}
  \item 赠与合同。\zd{双方、单务、无偿、诺成}。必考。
        \begin{enumerate}
          \item 赠与人。
                \begin{enumerate}
                  \item \zd{故意}不告知 \hz 保证无瑕疵 \yjt 瑕疵担保责任。\al{未告知马的习性。题目没说就没故意}。
                  \item 法定撤销权。\al{不叫法定解除权}。
                        \begin{enumerate}
                          \item 受赠人忘恩负义。严重侵犯人身（重伤）、有义务而不履行、附义务和附条件赠与。
                          \item 行权。有顺序。本人 1 年 \yjt 法代和继承人（6 个月）。
                        \end{enumerate}
                  \item 任意撤销权。
                        \begin{enumerate}
                          \item 例外。公证、公益、道德义务。\al{扶贫、助残、直系血亲之间、资助员工子女教育费用}。
                          \item vs 法定撤销权。用时间区分，转移财产前/后。
                        \end{enumerate}
                  \item 穷困抗辩权 \bdy 法定/任意撤销权。\al{罹患重病、花光继续、无法工作、生活陷入困难}。
                        \begin{enumerate}
                          \item 穷困抗辩是不再继续履行。而撤销是撤回全部，当作从来没有发生过，东西也得给我还回来。
                          \item 思路 2。法定撤销情形就 4 种。\zd{甭管题目描述得多天花乱坠}。
                        \end{enumerate}
                \end{enumerate}
          \item 附加的负担仅为合同上的限制，不得对抗第三人。
        \end{enumerate}
\end{enumerate}
